\section{Introduction}

Spectral localization in the standard Hilbert space of a von Neumann
algebra does not by itself produce uniformly bounded elements of the
algebra. This distinction matters when an argument starts with spectral
vectors but its asymptotic identities apply only to bounded operator
sequences. Our main analytical result, \Cref{rec:bounded}, provides a
bounded recovery principle for this situation.

Let $M$ be a von Neumann algebra with a faithful normal state $\phi$.
Write $M_\phi$ for the fixed-point algebra of its modular automorphism
group $\sigma^\phi$. We call $\phi$ ergodic when $M_\phi=\C1$.
On the standard Hilbert space $H=L^2(M)$, put
$\xi=\phi^{1/2}$, $D=\log\Delta_\phi$, and $U_t=e^{itD}$.
For an arbitrary fixed $T\in\B(H)$, consider unit vectors supported in
spectral intervals of $D$ shrinking to a real number $s$.
\Cref{rec:bounded} shows that a positive limiting averaged value of
$\norm{TU_th}^2$ on these vectors can be recovered as a positive value of
$\norm{Tv\xi}^2$ on a sequence of uniformly bounded $v\in M$.
The recovered vectors remain in intervals shrinking to $s$, with widths
at most four times the original widths. The averaging uses a fixed
translation-invariant mean obtained from symmetric interval averages.
The exact quantifiers and mean are given in
\Cref{mod:average,rec:bounded}.

The only structural assumption in this recovery theorem is
$M_\phi=\C1$; neither a bicentralizer condition nor a second action is
required. This assumption supplies a broad class of applications:
Marrakchi and Vaes proved that ergodic faithful normal states form a
dense $G_\delta$ in the normal state space of every type
$\mathrm{III}_1$ factor with separable predual
\cite[Theorem~A]{MarrakchiVaes2024}.

\paragraph{Bicentralizers and prior work.}
We use the predual convention
\[
  (w\phi)(y)=\phi(yw),\qquad
  (\phi w)(y)=\phi(wy).
\]
The bicentralizer $\BC(M,\phi)$ consists of those $a\in M$ for which
$[a,w_n]\to0$ strongly whenever $(w_n)$ is uniformly bounded and
$\norm{w_n\phi-\phi w_n}_{M_*}\to0$.
Connes' bicentralizer conjecture asserts that this algebra is scalar
for every type $\mathrm{III}_1$ factor with separable predual.
Connes developed the problem in his classification program for injective
factors \cite{Connes1980,Connes1985}. Haagerup proved the amenable case,
completing the uniqueness theorem for the injective type
$\mathrm{III}_1$ factor \cite{Haagerup1987}.

Haagerup also related triviality to the existence of a faithful normal
state with irreducible centralizer, meaning
$(M_\phi)'\cap M=\C1$ \cite[Theorem~3.1]{Haagerup1987}.
Together with Popa's work on maximal abelian subalgebras
\cite{Popa1981}, this connects the problem to expected subalgebras;
the expected-masa equivalence is stated in
\cite[Corollary~3.6]{HI2017}.
Important nonamenable cases include Houdayer's theorem for free Araki--Woods factors
\cite{Houdayer2009} and the semisolid factors treated by Houdayer and
Isono \cite[Theorem~3.7]{HI2017}.
An irreducible centralizer in these results is different from the scalar
centralizer assumed in our recovery theorem.

Houdayer and Isono proved that a nonscalar bicentralizer of a separable
type $\mathrm{III}_1$ factor is itself a type $\mathrm{III}_1$ factor
and equals its own bicentralizer for the restricted state
\cite[Proposition~3.4 and Theorem~3.5]{HI2017}.
Ando, Haagerup, Houdayer, and Marrakchi established the canonical relative
bicentralizer flow \cite[Theorem~A]{AHHM2020}, following the flow
construction of Haagerup and its independent discovery by Marrakchi.
Marrakchi proved that every nonidentity time of the absolute flow is
ergodic, without a fullness assumption
\cite[Theorem~D]{Marrakchi2020}. His subsequent work develops the
binormal-state and relative-core machinery
\cite{Marrakchi2025,Marrakchi2026}.

\paragraph{The absolute application and the recent announcement.}
Houdayer and Marrakchi have announced triviality of the bicentralizer
for every type $\mathrm{III}_1$ factor admitting a faithful normal
state, with no separability condition in the theorem statement
\cite[Theorem~G]{HoudayerMarrakchi2026}.
The application proved here is an alternative proof of the
separable-predual instance of that conclusion.

\begin{theorem}\label{thm:main}
Let $M$ be a type $\mathrm{III}_1$ factor with separable predual and let
$\phi$ be any faithful normal state on $M$. Then
\[
  \BC(M,\phi)=\C1.
\]
\end{theorem}

By \cite[Corollary~3.6]{HI2017}, \Cref{thm:main} also implies that every
type $\mathrm{III}_1$ factor with separable predual admits a maximal
abelian subalgebra with a faithful normal conditional expectation.
This is an existence statement for the factor; the expectation need not
preserve the scalar-centralizer state used in bounded recovery.

The independent analytical contribution is the bounded recovery theorem, which applies
to arbitrary $T$ without assuming a bicentralizer intertwining law.
Houdayer and Marrakchi obtain the same symmetric spectral relation used
below through multiplication isometries and a singular KMS state on the
continuous core
\cite[Section~8.4, Propositions~8.8--8.10]{HoudayerMarrakchi2026}.
Here modular averaging and bounded recovery lead to that relation through
two interval estimates.

Houdayer and Marrakchi explicitly disclose assistance from
\emph{GPT-5.6 Sol} and \emph{GPT-6 Astra}
\cite[p.~7]{HoudayerMarrakchi2026}. They credit Sol with counterexample
exploration and an initial strict-outerness argument, and Astra with
the singular tracial-state ingredient in their resonance strategy.
They describe checking and rewriting the proofs and retaining
responsibility for them. These are their stated contributions to that
work, whose relative tracial-action theory is also developed in the paper.

\paragraph{The rigidity statement.}
A continuous action on $M$ means a $u$-continuous action on its
automorphism group. Bounded recovery gives the following spectral proof
of rigidity.

\begin{theorem}[Spectral intertwining rigidity]\label{thm:rigidity}
Let $M$ be a von Neumann algebra with a faithful normal state $\phi$ such
that $M_\phi=\C1$. Let $b:\R\curvearrowright M$ be a continuous
$\phi$-preserving action. Suppose that, for every $s\in\R$, every uniformly
bounded sequence $(w_n)$ in $M$, and every $a\in M$,
\begin{equation}\label{eq:flow-property}
  \norm{\phi w_n-e^s w_n\phi}_{M_*}\longrightarrow0
  \quad\Longrightarrow\quad
  w_na-b_s(a)w_n\longrightarrow0
  \quad\text{strongly}.
\end{equation}
Then $b_s=\id_M$ for every $s\in\R$.
\end{theorem}

The conclusion also follows from the announced absolute theorem.
Indeed, at $s=0$, hypothesis \eqref{eq:flow-property} implies
$\BC(M,\phi)=M$. If $M\ne\C1$, its scalar centralizer makes it a
type $\mathrm{III}_1$ factor by
\cite[Lemma~2.1]{MarrakchiVaes2024}; the announced theorem would then
force $M=\C1$. Our proof instead derives \Cref{thm:rigidity} directly
from bounded recovery and uses it to prove \Cref{thm:main}.
For that application to a self-bicentralizing factor, the modular
centralizer is scalar, the bicentralizer flow satisfies
\eqref{eq:flow-property}, and that flow is ergodic. A trivial ergodic
action has scalar underlying algebra.

\paragraph{Recovery and the spectral argument.}
The recovery proof combines modular averaging with the fourth-moment
and polar-cutoff devices of Haagerup and Musat
\cite[Equation~(2.10), Lemma~2.6, and Remark~2.8]{HaagerupMusat2007}.
The strategy of recovering bounded witnesses by randomization,
fourth-moment control, and truncation appears in the tracial spectral-gap
argument of Ioana and Vaes \cite[pp.~127--130]{IoanaVaes2015}.
Related fourth-moment control underlies Marrakchi's binormal-state
implementation \cite[Section~2, Claim~2.2]{Marrakchi2025}.
Here modular orbit averages supply the moment bounds, and a final Fourier
filter restores the shrinking spectral bands.
The averaging time is chosen separately for each initial bounded
representative of a spectral vector; these representatives need not be
uniformly bounded.
Norm convergence in the predual controls a varying second-moment
operator, and independent circle phases give a uniform fourth-moment
bound. One fixed polar truncation threshold then suffices. Fourier
localization restores the spectral support without increasing operator
norms by a factor depending on the width of the spectral interval.
The Fourier operations belong to the standard spectral-subspace
framework for automorphism groups \cite{Arveson1974}; the bounds needed
here are proved explicitly.

For the spectral application, implement the state-preserving action by
$V_s(a\xi)=b_s(a)\xi$ and write $V_s=e^{isQ}$. These unitaries commute
with $U_t$, so $(D,Q)$ has a joint spectral support $\mathcal S\subset\R^2$.
The recovered bounded sequences let \eqref{eq:flow-property} force
averaged intertwining defects to vanish on bands shrinking to a fixed
modular energy. A spectral partition then gives an interval estimate
without requiring a rate uniform in that energy. Applying this estimate
twice, with the operators interchanged and the time reversed, gives
\[
  \exp\bigl(i(r\ell+qp)\bigr)=1
  \qquad ((r,p),(q,\ell)\in\mathcal S).
\]
This symmetric relation is the final rigidity mechanism. It forces the
joint support to be countable or to lie on a coordinate axis.
Scalarity of the modular centralizer excludes nonzero modular
eigenvalues, by the bounded-eigenoperator extraction and KMS argument of
\cite[Lemma~2.1]{MarrakchiVaes2024}, leaving only the axis on which the
flow generator vanishes.

\paragraph{Organization.}
\Cref{sec:bicentralizers} records the established bicentralizer inputs
and reduces \Cref{thm:main} to \Cref{thm:rigidity}.
\Cref{sec:modular} proves the modular and averaging preliminaries.
\Cref{sec:recovery} proves \Cref{rec:bounded}, and
\Cref{sec:rigidity} completes the spectral argument.
